% Converted from README.md
\paragraph{Primer: Sortiranje}\label{primer-sortiranje}

Program izvršava funkcije koje implementiraju različite algoritme sortiranja, na rastuće ili opadajuće sortiranom nizu ili nasumice generisanom nizu celih brojeva. U zavisnosti od izbora ulaznog niza neki algoritmi se pokazuju daleko bolje od očekivanog.

Izvorni kod je podeljen na \texttt{main.c}, \texttt{sort.c} i \texttt{sort.h}. Program generiše niz, pravi kopiju za svaki algoritam i zatim redom pokreće selection sort, insertion sort, bubble sort, shell sort, merge sort i quick sort.

Primer se može prevesti ručno:

\begin{minted}{bash}
$ gcc -O2 -g -fno-inline main.c sort.c -o sort.out
\end{minted}

Pokretanje sa nasumičnim nizom:

\begin{minted}{bash}
$ valgrind --tool=callgrind ./sort.out 10000
\end{minted}

Pokretanje sa rastuće i opadajuće sortiranim nizom:

\begin{minted}{bash}
$ valgrind --tool=callgrind ./sort.out -r 10000
$ valgrind --tool=callgrind ./sort.out -o 10000
\end{minted}

U izveštaju treba porediti funkcije \texttt{selection\_sort}, \texttt{insertion\_sort}, \texttt{bubble\_sort}, \texttt{shell\_sort}, \texttt{merge\_sort} i \texttt{quick\_sort}. Posebno je korisno uporediti kako se broj instrukcija menja kada je ulaz već sortiran ili sortiran opadajuće. Na tim ulazima neki algoritmi imaju najbolji slučaj, a neki najgori slučaj.

Ako je izveštaj prevelik, smanjiti dimenziju niza. Cilj primera nije apsolutno merenje vremena, već poređenje oblika troška između algoritama.
